\color{unidarkgreen}{\subsection{Дистанционная защита}}
\color{black}

\begin{enumerate}[label=\arabic{section}.\arabic{subsection}.\arabic*,  align=left, labelsep=4pt, leftmargin=0pt, itemindent=57pt]

\item
ИЭУ содержит шесть ступеней ДЗ от междуфазных замыканий и две ступени ДЗ от замыканий на землю. Дистанционная  защита  (ДЗ)  предназначена  для  селективного  отключения  поврежденного  элемента
электрической сети при всех видах  коротких  замыканий (КЗ),  реагируя  на снижение измеряемого сопротивления в защищаемой зоне.

\item
Уставки ДЗ представлены в таблице \ref{app_set:dz}.

\item 
ДЗ селективно срабатывает:
\begin{itemize}
    \item при всех видах КЗ; 
    \item при переходе внешнего КЗ во внутреннее; 
    \item при КЗ, сопровождающемся снижением измеряемого напряжения до нуля.
\end{itemize}

\item 
Обеспечивается несрабатывание защиты:
\begin{itemize} 
\item при качаниях и асинхронном ходе; 
\item при неисправностях цепей напряжения;
\item при внешних КЗ с насыщением трансформаторов тока (при соблюдении заявленных производителем требований к ТТ; 
\item при бросках тока намагничивания без КЗ;
\item при внешних КЗ, даже если измеряемое напряжение снижается до нуля.
\end{itemize}

\item Каждая из ступеней ДЗ от междуфазных КЗ включает в себя три реле сопротивления (РС), использующие для замера сопротивления  комплексные значения линейных напряжений и соответствующих им линейных
токов:
% ФОРМУЛЫ
\begin{align}
\vec Z_{AB} =  \frac{\vec U_A - \vec U_B}{\vec I_A - \vec I_B}\label{eq:groupZ1}, \\[10pt]
\vec Z_{BC} =  \frac{\vec U_B - \vec U_C}{\vec I_B - \vec I_C}\label{eq:groupZ2}, \\[10pt]
\vec Z_{CA} =  \frac{\vec U_C - \vec U_A}{\vec I_C - \vec I_A}\label{eq:groupZ3}, 
\end{align}
\begin{eqexpl}[20mm]
    \item{$ \vec U_{A}, \vec U_{B}, \vec U_{C} $} комплексные значения напряжений фаз А, В, С;
	\vspace{2pt}
    \item{$ \vec I_{A}, \vec I_{B}, \vec I_{C} $} комплексные значения токов фаз А, В, С.
\end{eqexpl}
\vspace{4pt}

\item 
Комплексные значения сопротивлений по контуру «фаза-земля» рассчитываются с учетом компенсации тока нулевой последовательности:
% ФОРМУЛЫ
\begin{align}
\vec Z_{A} =  \frac{\vec U_A }{\vec I_A + \vec K \cdot  3 \vec I_0}\label{eq:groupZ11}, \\[10pt]
\vec Z_{B} =  \frac{\vec U_B }{\vec I_B + \vec K \cdot  3 \vec I_0}\label{eq:groupZ12}, \\[10pt]
\vec Z_{C} =  \frac{\vec U_C }{\vec I_C + \vec K \cdot  3 \vec I_0}\label{eq:groupZ13}, 
\end{align}
\begin{eqexpl}[5mm]
    \item{$ 3 \vec I_{0} $}  комплексное  значение  утроенного  тока  нулевой  последовательности  защищаемого объекта (расчетное значение);
    \item{$ \vec K $}  комплексный коэффициент компенсации тока нулевой последовательности защищаемого объекта.
\end{eqexpl}
\vspace{4pt}

\item 
В ИЭУ коэффициент  $\vec K$  задается двумя уставками, \textbf{«KR - коэффициент компенсации тока 3I0 по активному сопротивлению»}  и  \textbf{«KХ  -  коэффициент  компенсации тока  3I0 по реактивному сопротивлению»}.  

\item 
Коэффициенты компенсации рассчитываются по следующей формуле:
% ФОРМУЛЫ
\begin{align}
\vec K = \frac{1}{3} \cdot \frac{\vec Z_{0 \textsl{ уд}} - \vec Z_{1 \textsl{ уд}}}{\vec Z_{1 \textsl{ уд}}}\label{eq:groupK1},
\end{align}

\begin{eqexpl}
\item{$ \vec Z_{0 \textsl{ уд}} $}  комплексное удельное сопротивление нулевой последовательности;
\item{$ \vec Z_{1 \textsl{ уд}} $}  комплексное удельное сопротивление прямой последовательности.
\end{eqexpl}
\vspace{4pt}

\item 
Характеристики  срабатывания  ненаправленного  РС  для  междуфазных  и  фазных  контуров  задаются независимо для каждой ступени уставками:
\begin{itemize} 
\item \textbf{«Rср ДЗ-ФФ(ФЗ)-1(2..6)»} – активное сопротивление срабатывания РС; 
\item \textbf{«Хср ДЗ-ФФ(ФЗ)-1(2..6)»} – реактивное сопротивление срабатывания РС;
\item \textbf{«Ф1 ДЗ-ФФ(ФЗ)-1(2..6)»} – угол линии электропередач; 
\end{itemize}

\item 
Для  предотвращения  неселективного  срабатывания  РС-ФЗ-1  и  РС-ФФ-1  при  внешних  повреждениях через  переходное  сопротивление  предусмотрена  возможность  отсечения  области  характеристики
срабатывания,  задаваемой  углом  $φ_4$.  Область  отсекается  отрезком,  исходящей  под  углом  $φ_4$,  из  точки  на верхней части характеристики, соответствующей углу линии электропередачи  $φ_1$. Угол $φ_4$ задается уставкой \textbf{«Ф4 ДЗ-ФФ(ФЗ)-1»}. 

\item 
Вид характеристики срабатывания ненаправленного РС представлен на рисунке \ref{pict_dz:char1}. 

\begin{figure}[H]
\centering
\includegraphics[width=0.4\textwidth,height=0.4\textheight,keepaspectratio]{\fbpath/07.02 РЗАТ/051.0041 ДЗ/img/_char1.pdf}
\captionof{figure}{ Характеристика срабатывания ненаправленного РС}\label{pict_dz:char1}
\end{figure}

\item 
Параметры реле сопротивления приведены в таблице \ref{desc_dz:tbl1}

\small
\setlength{\extrarowheight}{0.05cm} % высота ячеек в таблице
\begin{longtable}
{|>{\centering\arraybackslash}m{15.2cm}|>{\centering\arraybackslash}m{1.7cm}|}
\caption{Параметры РС\hfill\vspace{-0.5\baselineskip}}\label{desc_dz:tbl1}\\ 
\hhline{|-|-|}
\rowcolor{gray!30}
Наименование параметра & Значение  \\ 
\hhline{|-|-|}
\endfirsthead

\caption*{\hspace{3pt}\emph{Продолжение таблицы \ref{desc_dz:tbl1}\hfill\vspace{-0.5\baselineskip}}} \\ % сделано по ГОСТ 2.105 п.6.8.7
\hhline{|-|-|}
\rowcolor{gray!30}
Наименование параметра & Значение  \\ 
\hhline{|-|-|}
\endhead
\hhline{|-|-|}

\endfoot
\endlastfoot

\raggedright Минимальное  междуфазное  напряжение,  при  котором  обеспечиваются  точностные характеристики РС, В, не менее & \centering\arraybackslash 0,5  \\  \hline
\raggedright Ток точной работы во всем диапазоне уставок при работе на угле $φ_1$, \% от номинального тока, не более  & \centering\arraybackslash 10  \\  \hline
\raggedright Собственное  время  срабатывания  РС  при  работе  на  угле  максимальной  чувствительности, трехкратном токе точной работы и скачкообразном уменьшении напряжения  на входе РС от напряжения,  соответствующего  сопротивлению  на  зажимах  РС  не  менее  $1,2\cdot(X_{\textsl{уст}} /sinφ_1 )$ до напряжения соответствующего $0,6\cdot(X_{\textsl{уст}} /sinφ_1 )$ мс, не более   & \centering\arraybackslash 25  \\  \hline
\raggedright Собственное  время  возврата  РС  при  работе  на  угле  максимальной  чувствительности, трехкратном  токе  точной  работы  и  скачкообразном  увеличении  напряжения  на  входе  РС  от напряжения соответствующего сопротивлению на зажимах РС $0,1\cdot(X_{\textsl{уст}} /sinφ_1 )$ до напряжения соответствующего $1,2\cdot(X_{\textsl{уст}} /sinφ_1 )$ (но не более 100 В для междуфазного канала и 57 В –- для фазного), мс, не более   & \centering\arraybackslash 50  \\  \hline
\raggedright Средняя основная погрешность всех РС по величине сопротивления срабатывания $R_{\textsl{уст}}$  и $X_{\textsl{уст}}$ при токе, равном $I_{\textsl{ном}}$ (или, в зависимости от уставки, меньшем токе, исходя из максимального напряжения на зажимах РС, равного 100 В), \% от уставки, не более  & \centering\arraybackslash $\pm$5  \\  \hline
\raggedright Основная абсолютная погрешность РС по углу наклона характеристики при токе, равном $I_{\textsl{ном}}$, эл. град., не более   & \centering\arraybackslash $\pm$5  \\  \hline
\raggedright Дополнительная абсолютная погрешность РС по углу наклона характеристики при изменении тока КЗ в диапазоне $(0,2 \dots 30) \cdot I_{\textsl{ном}}$, эл. град. от значения, измеренного при $I_{\textsl{ном}}$, не более   & \centering\arraybackslash $\pm$7  \\  \hline
\raggedright Дополнительная  погрешность  РС  по  сопротивлениям  $R$  и  $X$  при  изменении  температуры окружающей  среды  в  рабочем  диапазоне,  \%  от  среднего  значения,  измеренного  при нормальных климатических условиях, не более   & \centering\arraybackslash $\pm$3  \\  \hline
\raggedright Время  срабатывания  ДЗ  при  переходе  внешнего  КЗ  во  внутреннее  в  условиях  наличия насыщения ТТ мс, не более & \centering\arraybackslash 60  \\  \hline
\raggedright Коэффициент возврата, не более   & \centering\arraybackslash 1,1  \\  \hline

\end{longtable}
\setlength{\extrarowheight}{0.0cm} % возврат параметра высота ячеек в таблице
\normalsize

\item 
Выполнение функций направленного РС обеспечивается совместным использованием ненаправленного РС и органов направленности (ОН). Орган направленности представляет собой два отрезка, исходящие из
начала координат и расположенные во втором и четвёртом квадрантах. Вид характеристики  (см. рис. \ref{pict_dz:on}) определяется углами $φ_2$ и $φ_3$  наклона  этих  отрезков,  отсчитываемых  относительно  оси  $R$.  Углы  наклона  характеристики задаются уставками \textbf{«Ф2 ОН-ФФ(ФЗ)»} и\textbf{«Ф3 ОН-ФФ(ФЗ)»}. Уставки органов направленности являются общими
для всех ступеней, задаются раздельно для контуров «фаза-фаза» и «фаза-земля». Углы наклона в ИЭУ задаются  для  прямонаправленного  ОН в сторону АТ,  а  характеристика  обратнонаправленного  ОН  симметрична характеристике прямонаправленного ОН с поворотом на 180 градусов.

\begin{figure}[H]
\centering
\includegraphics[width=0.6\textwidth,height=0.6\textheight,keepaspectratio]{\fbpath/07.02 РЗАТ/051.0041 ДЗ/img/_on.pdf}
\captionof{figure}{Характеристики срабатывания органов направления}\label{pict_dz:on}
\end{figure}

\item
Направленность для каждой ступени ДЗ задается с помощью программных переключателей \textbf{«Направл. ДЗ-ФФ(ФЗ)-1(2..6)»}, имеющих три положения:
\begin{enumerate}[label=\arabic*)]
\item \textit{Ненапр.} – ступень выполнена ненаправленной; 
\item \textit{В АТ} – ступень выполнена прямонаправленной в сторону АТ; 
\item \textit{В сеть} – ступень выполнена обратнонаправленной в сторону сети. 
\end{enumerate}

\item
При  близких  трехфазных  КЗ,  когда  все  напряжения  на  входе  РС  близки  к  нулю,  для  определения направленности в течение времени не менее 40 мс используются напряжения предаварийного режима (работа по «памяти»). При работе РС «по памяти» при трехфазных КЗ в месте установки защиты обеспечивается длительность сигнала срабатывания на выходе РС не менее 40 мс в диапазоне токов от $0,2\cdot I_{\textsl{ном}}$ до $30\cdot I_{\textsl{ном}}$. При КЗ «за спиной» с токами до $30\cdot I_{\textsl{ном}}$  обеспечивается отсутствие ложных срабатываний РС, в том числе при снижении измеряемого напряжения до нуля. 

\item
Для отстройки от нагрузочного режима возможен вырез в характеристике срабатывания РС. Предусмотрено два органа выреза нагрузочного режима для ступеней направленных в АТ и для ступеней направленных в сеть, уставки органов выреза в АТ и в сеть для контуров «фаза-фаза» и «фаза-земля» задаются раздельно.  Характеристика  срабатывания  приведена  на  рисунке \ref{pict_dz:r}.  Эта характеристика срабатывания определяется уставками: 
\begin{itemize}
    \item \textbf{«Rнг-ФФ(ФЗ)в АТ»} - сопротивление выреза нагрузочного режима ступеней, направленных в сторону АТ;
    \item \textbf{«Фнг-ФФ(ФЗ) в АТ»} - угол выреза нагрузочного режима ступеней, направленных в сторону АТ;
	\item \textbf{«Rнг-ФФ(ФЗ) в сеть»} - сопротивление выреза нагрузочного режима ступеней, направленных в сторону сети;
	\item \textbf{«Фнг-ФФ(ФЗ) в сеть»} - угол выреза нагрузочного режима ступеней, направленных в сторону сети.
\end{itemize}

\item
Исключаемая область в характеристике срабатывания симметрична относительно осей координат. Ввод в работу отстройки от нагрузки в сторону АТ и в сторону сети осуществляется с помощью программных
переключателей \textbf{«Zнг в АТ»} и \textbf{«Zнг в сеть»}  соответственно. 

\begin{figure}[H]
\centering
\includegraphics[width=0.5\textwidth,height=0.5\textheight,keepaspectratio]{\fbpath/07.02 РЗАТ/051.0041 ДЗ/img/_r.pdf}
\captionof{figure}{Характеристика срабатывания измерительного органа отстройки от нагрузки}\label{pict_dz:r}
\end{figure}

\item
Суммарная характеристика направленного РС, используемого в логике работы ступеней ДЗ приведена на рисунке \ref{pict_dz:rs}.  

\begin{figure}[H]
\centering
\includegraphics[width=0.75\textwidth,height=0.75\textheight,keepaspectratio]{\fbpath/07.02 РЗАТ/051.0041 ДЗ/img/_rs.pdf}
\captionof{figure}{Характеристики срабатывания направленного РС}\label{pict_dz:rs}
\end{figure}

%%%%%%%%%%%%%%%%%%%%%%%%%%%%%%%%%%%% ООВП  %%%%%%%%%%%%%%%%%%%%%%%%%%%%%%%%%%

\item
Для  обнаружения  однофазных  КЗ  и  пуска  ступеней  ДЗ-ФЗ,  контролирующих  контуры  «фаза-земля», предусмотрен  орган  определения  вида  повреждения  (ООВП),  основанный  на  совместном  применении
измерительного органа тока нулевой последовательности (ИО 3I0 ООВП) с торможением от одного из фазных токов (средний из  $I_A , I_B , I_C$) и измерительного органа напряжения нулевой последовательности (ИО 3U0 ООВП). Торможение осуществляется от тока той фазы, значение в которой является средним между максимальным и минимальным значениями токов в остальных двух фазах. В таком случае при междуфазном КЗ на землю торможение максимальном, а при однофазном КЗ –- минимально. 

\item
Характеристика  срабатывания  ИО  3I0  ООВП  с  торможением  представлена  на  рисунке  \ref{pict_dz:torm}.  

\begin{figure}[H]
\centering
\includegraphics[width=0.65\textwidth,height=0.65\textheight,keepaspectratio]{\fbpath/07.02 РЗАТ/051.0041 ДЗ/img/_torm.pdf}
\captionof{figure}{Характеристика срабатывания ИО 3I0 ООВП}\label{pict_dz:torm}
\end{figure}

\item
Порог срабатывания горизонтального участка задается уставкой \textbf{«3I0ср ООВП»}. Наклонный участок характеристики проходит через точку, равной $1,25 \cdot I_{\textsl{ном}}$ , под углом к горизонтальной оси, соответствующим коэффициенту торможения $k_T$. Коэффициент торможения задается уставкой \textbf{«Кт ООВП»}. При превышении тормозного тока уставки \textbf{«Iср блок. торм.»} ООВП блокируется. Алгоритм работы логики ООВП приведен на рисунке \ref{pict_dz:OOVP}.

\medskip
\begin{figure}[H]
\centering
\includegraphics[width=0.8\textwidth,height=0.8\textheight,keepaspectratio]{\fbpath/07.02 РЗАТ/051.0041 ДЗ/img/OOVP.pdf}
\captionof{figure}{Функциональная схема логики ООВП}\label{pict_dz:OOVP}
\end{figure}

%%%%%%%%%%%%%%%%%%%%%%%%%%%%% СТУПЕНИ ДЗ %%%%%%%%%%%%%%%%%%%%%%%%%%%%%%%%%%%%%

\item
Ввод  в  работу  ступеней  ДЗ,  работающих  по  контуру  «фаза-земля»  может  осуществляться  с дополнительным  контролем  напряжения  нулевой  последовательности  за  счет  ввода  программного
ключа \textbf{«Контр.  3U0 ООВП»}.  Напряжение  срабатывания  ИО  3U0  ООВП  задается  уставкой  \textbf{«3U0ср ООВП»}. ИО 3U0 ООВП может работать по напряжению $3U_{\textsl{0 звезда}}$, рассчитанному на основе фазных напряжений, или по напряжению $3U_{\textsl{0 треуг.}}$   --  выбирается  программным  ключом  \textbf{«Выбор  3U0»}, имеющим  два  соответствующих положения: \textsl{«Расчетн.»} и \textsl{«Измерен.»}.

\item
Переменная $3U_{\textsl{0 треуг.}}$ зависит от положения программного переключателя \textbf{«Схема подключения цепей разомкнутого треугольника»}.  Данный  переключатель  позволяет  учесть  различные  схемы  соединения дополнительной  обмотки  ТН.  Если  переключатель  выставлен  в  положение  «$U_{\textsl{НК}}$»,  то  переменной  $3U_{\textsl{0 треуг.}}$ присваивается  значение  напряжения,  измеренного  по  аналоговому  входу  «$U_{\textsl{ИФ}}(U_{\textsl{НК}})$».  Если  программный переключатель {\textbf{«Схема подключения цепей треугольника»}} выставлен в положение «$U_{\textsl{НИ}}/U_{\textsl{ИК}}$», «$U_{\textsl{НФ}}/U_{\textsl{ФК}}$» или  «$U_{\textsl{НИ}}/U_{\textsl{ИФ}}/U_{\textsl{ФК}}$», то переменная $3U_{\textsl{0 треуг.}}$ рассчитывается в соответствии с формулами, приведенными в приложении \ref{app:signals}. Программный переключатель \textbf{«Выбор 3U0»} аналогично влияет на работу ОНМНП. 

\item
В составе логики дистанционной защиты предусмотрено шесть ступеней от междуфазных замыканий с контролем трех контуров «фаза-фаза». 

\vspace{2mm}
\begin{figure}[H]
\centering
\includegraphics[width=0.8\textwidth,height=0.8\textheight,keepaspectratio]{\fbpath/07.02 РЗАТ/051.0041 ДЗ/img/DZ_FF_1.pdf}
\captionof{figure}{Функциональная схема логики ДЗ-ФФ-1}\label{pict_dz:DZ_FF_1}
\end{figure}

\item
Пороги срабатывания РС ступеней ДЗ-ФФ задаются уставками \textbf{«Хср-ФФ-1(2…6)»}, \textbf{«Rср-ФФ-1(2…6)»} и \textbf{«Ф1-ФФ-1(2…6)»}, а выдержки времени срабатывания ступеней регулируются уставками \textbf{«Тср ДЗ-ФФ-1(2…6)»}. Ввод в работу ступеней ДЗ-ФФ осуществляется с помощью программных переключателей \textbf{«ДЗ-ФФ-1(2…6)»}. Направленности для каждой ступени задается с помощью программных переключателей \textbf{«Направл ДЗ-ФФ-1(2…6)»}.

\item
Предусмотрена возможность ввода подхвата направленного пуска ступеней ДЗ-ФФ от ненаправленного пуска той же ступени, задается программным переключателем \textbf{«Подхват РС-ФФ-1(2…6)»}. Тем самым
гарантируется правильное действие защиты при близких замыканиях. 

\item
Оперативный ввод/вывод ДЗ-ФФ производится с помощью БУ (блоков управления функциями) \textit{«Вывод ДЗ»}.

\item
Пуски степеней ДЗ-ФФ направленные в сеть, действуют в логику предварительного деления своей стороны АТ. С помощью выдержек времени в логике деления \textbf{«Тср ДЗ-ФФ-1(2…6) на СВ»}, \textbf{«ДЗ-ФФ-1(2…6) на В сеч»} и \textbf{«Тср ДЗ-ФФ-1(2…6) на В»}, первым действием формируются сигналы на разделение секции (отключение СВ) или деление по сечению (отключение В сечения), вторым действием формируются сигналы на отключение выключателей своей стороны с пуском УРОВ и запретом АПВ. Пуски ступеней ДЗ-ФФ направленные в АТ, могут действовать в логику предварительного деления резервных защит смежной стороны (смежный комплект РЗАТ), через выдержки времени в логике деления \textbf{«Тср ДЗ-ФФ-1(2…6) на пуск ЛД см ст»}. Действие срабатывания ступеней ДЗ-ФФ предусматривает отключение АТ со всех сторон с запретом АПВ и пуском УРОВ.

\item
В составе логики дистанционной защиты предусмотрено две ступеней от однофазных замыканий с контролем трех контуров «фаза-земля».

\vspace{2mm}
\begin{figure}[H]
\centering
\includegraphics[width=0.8\textwidth,height=0.8\textheight,keepaspectratio]{\fbpath/07.02 РЗАТ/051.0041 ДЗ/img/DZ_FZ_1.pdf}
\captionof{figure}{Функциональная схема логики ДЗ-ФЗ-1}\label{pict_dz:DZ_FZ_1}
\end{figure}

\item
Пуск ступеней ДЗ-ФЗ разрешается только при срабатывании органа определения вида повреждения \textit{«Ввод ДЗ-ФЗ»}. Пороги срабатывания РС ступеней ДЗ-ФЗ задаются уставками \textbf{«Хср-ФЗ-1(2)»}, \textbf{«Rср-ФЗ-1(2)»} и \textbf{«Ф1-ФЗ-1(2)»}, а выдержки времени срабатывания ступеней регулируются уставками \textbf{«Тср ДЗ-ФЗ-1(2)»}. Ввод в работу ступеней ДЗ-ФЗ осуществляется с помощью программных переключателей \textbf{«ДЗ-ФЗ-1(2)»}. Направленности для каждой ступени задается с помощью программных переключателей \textbf{«Направл ДЗ-ФФ-1(2)»}.

\item
Предусмотрена возможность ввода подхвата направленного пуска ступеней ДЗ-ФЗ от ненаправленного пуска той же ступени, задается программными переключателями \textbf{«Подхват РС-ФЗ-1(2)»}. Тем самым
гарантируется правильное действие защиты при близких замыканиях. 

\item
Оперативный ввод/вывод ДЗ-ФЗ производится с помощью БУ (блока управления функциями) \textit{«Вывод ДЗ»}.

\item
Пуски степеней ДЗ-ФЗ направленные в сеть, действуют в логику предварительного деления своей стороны АТ. С помощью выдержек времени в логике деления \textbf{«Тср ДЗ-ФЗ-1(2) на СВ»}, \textbf{«ДЗ-ФФ-1(2) на В сеч»} и \textbf{«Тср ДЗ-ФФ-1(2) на В»}, первым действием формируются сигналы на разделение секции (отключение СВ) или деление по сечению (отключение В сечения), вторым действием формируются сигналы на отключение выключателей своей стороны с пуском УРОВ и запретом АПВ. Действие срабатывания ступеней ДЗ-ФЗ предусматривает отключение АТ со всех сторон с запретом АПВ и пуском УРОВ.

\item
Для исключения ложного действия при возникновении неисправности в цепях напряжения ступени ДЗ могут быть заблокированы от БНН. Блокировка от БНН вводится отдельно для каждой ступени ДЗ с помощью программных переключателей \textbf{«БНН ДЗ-ФФ-1(2…6)»} и \textbf{«БНН ДЗ-ФЗ-1(2)»}.

\item
Дистанционная защита контролируется блокировкой при качаниях, предотвращающей ложное срабатывание ДЗ при качаниях и асинхронном ходе. В устройстве реализовано два способа блокировки:
\begin{itemize}
\item по приращению токов прямой и обратной последовательности;
\item по скорости изменения величины сопротивления dZ/dt.
\end{itemize}

\item
Контроль ступеней ДЗ от блокировки при качаниях вводится программными переключателями \textbf{«БК ДЗ-ФФ(ФЗ)-1(2…6)»}, имеющими четыре положения:
\begin{itemize}
\item \textit{«Выведено»} -- ступень работает без контроля от БК;
\item \textit{«БК-I-б»} -- ступень пускается от сигнала ввода быстродействующих ступеней;
\item \textit{«БК-I-м»} -- ступень пускается от сигнала ввода медленнодействующих ступеней;
\item \textit{«БК-Z»} -- ступень пускается от сигнала БК по скорости изменения величины сопротивления.
\end{itemize}

\item
Если за время пуска от БК происходит срабатывание ненаправленного РС ступени ДЗ, то разрещение от БК подхватывается и удерживается до пропадания условий срабатывания ненаправленного РС соответствующей ступени.

\item
В составе логики дистанционной защиты предусмотрено четыре режима оперативного ускорения ДЗ.

\vspace{2mm}
\begin{figure}[H]
\centering
\includegraphics[width=0.85\textwidth,height=0.85\textheight,keepaspectratio]{\fbpath/07.02 РЗАТ/051.0041 ДЗ/img/OA_DZ.pdf}
\captionof{figure}{Функциональная схема логики ОУ ДЗ}\label{pict_dz:OA_DZ}
\end{figure}

\item
Оперативное ускорение ДЗ в сеть без выдержки времени. Предназначено для оперативного ускорения ступени ДЗ-ФФ при выполнении переключений в сети с действием на отключение всего АТ (с пуском УРОВ и
запретом АПВ) с одной выдержкой времени, которая регулируется уставкой \textbf{«Тср ОУ ДЗ»}, выбор ускоряемой ступени определяется программным переключателем \textbf{«ОУ ДЗ»}. Оперативный ввод ОУ ДЗ в сеть без выдержки времени производится с помощью команды управления \textit{«Ввод ОУ ДЗ в сеть без вв»}.

\item
Оперативное ускорение ДЗ в сеть с первой выдержкой времени. Предназначено для оперативного ускорения ступени ДЗ-ФФ при выводе ДЗО с действием на отключение всего АТ (с пуском УРОВ и запретом АПВ) с одной выдержкой времени, которая регулируется уставкой \textbf{«Тср ОУ1 ДЗ»}, выбор ускоряемой ступени определяется программным переключателем \textbf{«ОУ1 ДЗ»}. Оперативный ввод ОУ ДЗ в сеть с первой выдержкой времени производится с помощью команды управления \textit{«Ввод ОУ ДЗ в сеть с вв1»}.

\item
Оперативное ускорение ДЗ в сеть со второй выдержкой времени. Предназначено для оперативного ускорения ступени ДЗ-ФФ при выводе ДЗШ с действием пуска ОУ2 ДЗ в логику предварительного деления своей стороны. Действие срабатывания ОУ2 ДЗ предусматривает отключение всего АТ (с пуском УРОВ и запретом АПВ). Время срабатывания ОУ2 ДЗ регулируется уставкой \textbf{«Тср ОУ2 ДЗ»}, а выбор ускоряемой
ступени определяется программным переключателем \textbf{«ОУ2 ДЗ»}. Оперативный ввод ОУ ДЗ в сеть со второй выдержкой времени производится с помощью команды управления \textit{«Ввод ОУ ДЗ в сеть с вв2»}.

\item
Оперативное ускорение ДЗ в АТ. Предназначено для оперативного ускорения ступени ДЗ-ФФ при выводе основных защит АТ с действием на отключение всего АТ (с пуском УРОВ и запретом АПВ) с одной выдержкой времени, которая регулируется уставкой \textbf{«Тср ОУ ДЗ в АТ»}, выбор ускоряемой ступени определяется программным переключателем \textbf{«ОУ ДЗ в АТ»}. Оперативный ввод ОУ ДЗ в АТ производится с
помощью команды управления \textit{«Ввод ОУ ДЗ в АТ»}.

\item
В составе логики дистанционной защиты предусмотрен пуск автоматического ускорения ступеней ДЗ.

\vspace{2mm}
\begin{figure}[H]
\centering
\includegraphics[width=1\textwidth,height=1\textheight,keepaspectratio]{\fbpath/07.02 РЗАТ/051.0041 ДЗ/img/AA_DZ.pdf}
\captionof{figure}{Функциональная схема АУ ДЗ}\label{pict_dz:AA_DZ}
\end{figure}

\item
С помощью программного переключателя \textbf{«Пуск АУ ДЗ в АТ»} имеется возможность выбрать от какой ступени ДЗ от междуфазных коротких замыканий будет пускаться автоматическое ускорение при опробовании
АТ.

\item
С помощью программного переключателя \textbf{«Пуск АУ ДЗ в сеть»} имеется возможность выбрать от какой ступени ДЗ от междуфазных коротких замыканий будет пускаться автоматическое ускорение при опробовании ошиновки АТ или системы шин.

\item
Дополнительно предусмотрена возможность пуска ускорения от ненаправленных РС ступеней ДЗ от междуфазных коротких замыканий c контролем цепей напряжения (БНН) и пуском от БК. С помощью программного переключателя \textbf{«Пуск АУ ДЗ ненапр»} имеется возможность выбрать от РС какой ступени будет пускаться автоматическое ускорение при включении выключателя. Контроль пуска АУ ненаправленного РС
ступени от БК задается уставкой \textbf{«БК АУ ДЗ ненапр»}, а так же дополнительный контроль от неисправности цепей напряжения задается уставкой \textbf{«БНН АУ ДЗ ненапр»}.

\item
В составе логики дистанционной защиты предусмотрен режим ускорения от защит смежной стороны.

\vspace{2mm}
\begin{figure}[H]
\centering
\includegraphics[width=0.8\textwidth,height=0.8\textheight,keepaspectratio]{\fbpath/07.02 РЗАТ/051.0041 ДЗ/img/TA_DZ.pdf}
\captionof{figure}{Функциональная схема логики ускорения ДЗ от защит смежной стороны}\label{pict_dz:TA_DZ}
\end{figure}

\item
Предусматривается возможность встречного ускорения ДЗ от сигнала пуска ступени ДЗ \textit{«Уск. от ДЗ см.
ст»}, смежного комплекта резервных защит автотрансформатора (РЗАТ). Так же, с помощью программного
переключателя \textbf{«Уск ДЗ от РПО см ст»} предусматривается возможность разрешения встречного ускорения
ДЗ при наличии сигнала отключенного положения выключателей смежной стороны \textit{«См. ст. отключена»}.
Данный вид ускорения позволяет использовать ДЗ как защиту абсолютной селективности, так как пуски защит
направленных в АТ на комплектах РЗАТ ВН и СН указывают на КЗ внутри АТ. С помощью программного
переключателя \textbf{«Уск см ДЗ»} имеется возможность выбрать пуском какой ступени ДЗ-ФФ будет
контролироваться ускорение от защит смежной стороны, а также формироваться сигнал пуска ускорения в
защиты смежной стороны \textit{«Уск. ДЗ на см. ст»}. Действие защиты предусматривает отключение АТ со всех
сторон с пуском УРОВ и запретом АПВ, минуя этапы деления. Выдержка времени срабатывания ускорения ДЗ
от защит смежной стороны регулируется уставкой \textbf{«Тср уск ДЗ от см ст»}.



\end{enumerate}
